\section*{Capítulo \ref{ch:g12:proton-nmr} --- RMN do próton}

\begin{solution}{exo:g12:proton-nmr:1}
Metano: 1. Etano: 1 (os dois \ce{CH3} são iguais). Propano: 2 (dois
\ce{CH3} iguais, um \ce{CH2}). Metanol: 2 (\ce{CH3} e \ce{OH}).
Metoximetano: 1.
\end{solution}

\begin{solution}{exo:g12:proton-nmr:2}
2 vizinhos: um tripleto, 1:2:1. 3 vizinhos: um quarteto, 1:3:3:1.
0 vizinhos: um singleto.
\end{solution}

\begin{solution}{exo:g12:proton-nmr:3}
Total de \qty{40}{mm} para 8 \omterm{def:g9:inside-the-atom:proton}{prótons}: \qty{5}{mm} por \omterm{def:g9:inside-the-atom:proton}{próton}. Os sinais
representam 3, 2 e 3 \omterm{def:g9:inside-the-atom:proton}{prótons}.
\end{solution}

\begin{solution}{exo:g12:proton-nmr:4}
\omterm{def:g11:functional-groups:carbonyl}{Aldeído}: de 9.7 a \qty{10.0}{ppm}. Um \ce{CH3} de cadeia: de 0.7 a
\qty{1.3}{ppm}.
\end{solution}

\begin{solution}{exo:g12:proton-nmr:5}
Os seus doze \omterm{def:g12:proton-nmr:equivalent}{prótons equivalentes} dão uma única linha, fina e intensa,
fácil de fixar no zero. O \omterm{def:g6:solutions-and-solubility:solute}{solvente} deuterado não dá nenhum sinal de
\omterm{def:g9:inside-the-atom:proton}{próton} próprio, que senão encobriria os dos poucos miligramas de
amostra.
\end{solution}

\begin{solution}{exo:g12:proton-nmr:6}
Dois sinais. \ce{CH2} vizinho do cloro: 2 \omterm{def:g9:inside-the-atom:proton}{prótons}, um quarteto (3
vizinhos), na faixa 2.5--\qty{4.0}{ppm}. \ce{CH3}: 3 \omterm{def:g9:inside-the-atom:proton}{prótons}, um
tripleto (2 vizinhos), perto da faixa usual do \ce{CH3}, um pouco acima
porque o cloro está a dois \omterm{def:g7:atoms-and-molecules:atom}{átomos} de distância.
\end{solution}

\begin{solution}{exo:g12:proton-nmr:7}
Três sinais: os dois \ce{CH3}, 6 \omterm{def:g12:proton-nmr:equivalent}{prótons equivalentes}, um dupleto (1
vizinho, o \ce{CH}); o \ce{CH}, 1 \omterm{def:g9:inside-the-atom:proton}{próton}, desdobrado por 6 vizinhos em
7 linhas, na faixa do \ce{H-C-O} (3.3--\qty{4.5}{ppm}); o \ce{OH}, 1
\omterm{def:g9:inside-the-atom:proton}{próton}, um singleto.
\end{solution}

\begin{solution}{exo:g12:proton-nmr:8}
(a) propanona; (b) etanoato de etila; (c) ácido propanoico.
\end{solution}

\begin{solution}{exo:g12:proton-nmr:9}
O \ce{CH2} está ligado ao \omterm{def:g7:atoms-and-molecules:atom}{átomo} de oxigênio do \omterm{def:g11:functional-groups:carboxylic-acid}{éster}, que é
eletronegativo e desblinda os seus \omterm{def:g9:inside-the-atom:proton}{prótons}: faixa do \ce{H-C-O},
3.3--\qty{4.5}{ppm}. O \ce{CH3} está um carbono mais longe e fica na
faixa das cadeias.
\end{solution}

\begin{solution}{exo:g12:proton-nmr:10}
Propanona: um \ce{C=O} (\omterm{def:g11:functional-groups:carbonyl}{cetona}, \qty{1715}{cm^{-1}}) e seis \omterm{def:g12:proton-nmr:equivalent}{prótons
equivalentes}, um singleto. O propanal daria três sinais, um deles, o do
\omterm{def:g9:inside-the-atom:proton}{próton} aldeídico de \ce{CH3-CH2-CHO}, perto de 9.7--\qty{10.0}{ppm}.
\end{solution}

\begin{solution}{exo:g12:proton-nmr:11}
O \omterm{def:g9:inside-the-atom:proton}{próton} do \ce{OH} é trocado rapidamente entre as \omterm{def:g7:atoms-and-molecules:molecule}{moléculas}: ele dá um
singleto e não desdobra os seus vizinhos. O \ce{CH2} é portanto
desdobrado só pelos 3 \omterm{def:g9:inside-the-atom:proton}{prótons} do \ce{CH3}: um quarteto.
\end{solution}

\begin{solution}{exo:g12:proton-nmr:12}
Ácido butanoico \ce{CH3-CH2-CH2-COOH}: 4 sinais, um tripleto (\ce{CH3}),
um sinal de 6 linhas (\ce{CH2} central, 5 vizinhos), um tripleto
(\ce{CH2-C=O}), um singleto perto de 11--\qty{12}{ppm}. Etanoato de
etila: um singleto, um quarteto, um tripleto. Propanoato de metila \ce{CH3-CH2-CO-O-CH3}:
um tripleto, um quarteto (\ce{CH2-C=O}, 2.0--\qty{2.4}{ppm}), um
singleto (\ce{O-CH3}). Metanoato de propila \ce{H-CO-O-CH2-CH2-CH3}: 4
sinais, um singleto bem à esquerda (o H do carbono do \ce{C=O}), um
tripleto (\ce{O-CH2}), um sinal de 6 linhas, um tripleto (\ce{CH3}).
\end{solution}

\begin{solution}{exo:g12:proton-nmr:13}
O sinal do \ce{OH}: os \omterm{def:g9:inside-the-atom:proton}{prótons} de \ce{O-H} são trocados com o deutério
da água pesada, e o \ce{O-D} não dá sinal de \omterm{def:g9:inside-the-atom:proton}{próton}. Um sinal que some
depois da agitação com \ce{D2O} pertence a um \omterm{def:g9:inside-the-atom:proton}{próton} ligado a O ou a N.
\end{solution}

\begin{solution}{exo:g12:proton-nmr:14}
O etanol acrescenta um quarteto perto de \qty{3.69}{ppm}, um tripleto
perto de \qty{1.23}{ppm} e um singleto do \ce{OH}. O singleto do \omterm{def:g11:functional-groups:carboxylic-acid}{éster}:
\qty{30}{mm} para 3 \omterm{def:g9:inside-the-atom:proton}{prótons}, \qty{10}{mm} por \omterm{def:g9:inside-the-atom:proton}{próton} de \omterm{def:g11:functional-groups:carboxylic-acid}{éster}. O
tripleto do etanol: \qty{3}{mm} para 3 \omterm{def:g9:inside-the-atom:proton}{prótons}, \qty{1}{mm} por \omterm{def:g9:inside-the-atom:proton}{próton}
de etanol. Como cada \omterm{def:g7:atoms-and-molecules:molecule}{molécula} contribui com um \ce{CH3} para esses
sinais, as quantidades estão na razão $1/10$: cerca de uma \omterm{def:g7:atoms-and-molecules:molecule}{molécula} de
etanol para dez de \omterm{def:g11:functional-groups:carboxylic-acid}{éster}.
\end{solution}

\begin{solution}{exo:g12:proton-nmr:15}
Quatro sinais: os dois \ce{CH3} (6 \omterm{def:g9:inside-the-atom:proton}{prótons}), um dupleto; o \ce{CH} (1
\omterm{def:g9:inside-the-atom:proton}{próton}), desdobrado por 6 + 2 = 8 vizinhos em 9 linhas; o \ce{CH2} (2
\omterm{def:g9:inside-the-atom:proton}{prótons}), vizinho do \ce{O}, um dupleto, na faixa do \ce{H-C-O}; o
\ce{OH} (1 \omterm{def:g9:inside-the-atom:proton}{próton}), um singleto. O sinal do \ce{CH} é o que tem mais linhas.
\end{solution}

\begin{solution}{pb:g12:proton-nmr:1}
\textbf{1.} Um \omterm{def:g11:organic-skeletons:alkane}{alcano} de quatro carbonos é \ce{C4H10}; um \ce{C=O} tira
dois hidrogênios e os dois oxigênios não acrescentam nenhum:
\ce{C4H8O2}, a fórmula dos ácidos e dos \omterm{def:g11:functional-groups:carboxylic-acid}{ésteres} de quatro carbonos.

\textbf{2.} \ce{CH3-CH2-CH2-COOH} é o ácido butanoico;
\ce{CH3-CO-O-CH2-CH3}, o etanoato de etila; \ce{CH3-CH2-CO-O-CH3}, o
propanoato de metila; e \ce{H-CO-O-CH2-CH2-CH3}, o metanoato de
propila.

\textbf{3.} Os três \omterm{def:g11:functional-groups:carboxylic-acid}{ésteres} compartilham o grupo \omterm{def:g11:functional-groups:carboxylic-acid}{éster}; os quatro têm um
\ce{C=O}.

\textbf{4.} Um \ce{C=O}, na faixa dos \omterm{def:g11:functional-groups:carboxylic-acid}{ésteres} (perto de
\qty{1735}{cm^{-1}}).

\textbf{5.} A banda \ce{O-H} muito larga de 2500 a
\qty{3300}{cm^{-1}}: ausente, logo não é o ácido butanoico.

\textbf{6.} Não: os três \omterm{def:g11:functional-groups:carboxylic-acid}{ésteres} têm os mesmos grupos.

\textbf{7.} Três.

\textbf{8.} Um grupo etila \ce{CH3-CH2-}: o \ce{CH2} desdobrado por 3
\omterm{def:g9:inside-the-atom:proton}{prótons} (quarteto), o \ce{CH3} por 2 (tripleto).

\textbf{9.} Os seus \omterm{def:g9:inside-the-atom:proton}{prótons} não têm \omterm{def:g9:inside-the-atom:proton}{prótons} nos \omterm{def:g7:atoms-and-molecules:atom}{átomos} vizinhos: um
\ce{CH3} ligado ao carbono do \ce{C=O} ou ao oxigênio do \omterm{def:g11:functional-groups:carboxylic-acid}{éster}.

\textbf{10.} O \ce{CH2} está ligado ao oxigênio (faixa do \ce{H-C-O}).

\textbf{11.} Um tripleto (\ce{CH3}), um quarteto perto de
2.0--\qty{2.4}{ppm} (\ce{CH2} vizinho do \ce{C=O}) e um singleto na
faixa do \ce{H-C-O} (\ce{O-CH3}). O quarteto ficaria bem abaixo de
\qty{4.12}{ppm} e o singleto bem acima de \qty{2.04}{ppm}: não é o
desconhecido.

\textbf{12.} Quatro sinais (um singleto bem à esquerda, um tripleto, um
sinal de 6 linhas, um tripleto): não é o desconhecido, que tem três.

\textbf{13.} Etanoato de etila, \ce{CH3-CO-O-CH2-CH3}: singleto
\qty{2.04}{ppm} = \ce{CH3-C=O}; quarteto \qty{4.12}{ppm} = \ce{O-CH2};
tripleto \qty{1.26}{ppm} = o \ce{CH3} da ponta.

\textbf{14.} Singleto 3, quarteto 2, tripleto 3.

\textbf{15.} $72 / 8 = \qty{9}{mm}$ por \omterm{def:g9:inside-the-atom:proton}{próton}.

\textbf{16.} Singleto $3 \times 9 = \qty{27}{mm}$; tripleto
\qty{27}{mm}.

\textbf{17.} $18 + 27 + 27 = \qty{72}{mm}$.

\textbf{18.} Sim: os \omterm{def:g11:functional-groups:carboxylic-acid}{ésteres} pequenos são conhecidos pelos seus cheiros
de fruta.

\textbf{19.} O degrau do quarteto mede \qty{18}{mm} dos \qty{72}{mm}.
\end{solution}
